% AINTEC 2026 Poster
% Author list is real names (user call 2026-09-21) — unlike paper.tex's \author{Anonymous}.
%
% 2026-09-24: operator / ISP market-structure material removed on อาจารย์'s instruction —
% it is now a separate poster owned by Kunanont Malayanont. Content bank for it:
% docs/poster/isp/content.md. Nothing operator-layer belongs on this poster any more;
% that includes the "operators grouped by ASN" method bullet and the 22-39 Mbps
% leading-mobile-operator band, both of which moved with it.
%
% Number provenance: body text quotes docs/paper/paper.tex; figures marked "paper" are
% docs/paper/figures/imageN.png (identical to the paper). Figures 01 and 03 come from
% scripts/build_cross_country_figs_8c.py, regenerated 2026-09-21 because that script
% silently excluded Singapore. Chart 03 (capital vs national) gives Singapore capital
% 361.0 vs paper prose 341; ratio 16x (361.0/22.0) vs paper prose 12x. Text next to the
% chart quotes the chart. Confirm with paper author which is current.
\documentclass[25pt, a0paper, portrait, margin=0mm, innermargin=15mm,
  blockverticalspace=9mm, colspace=15mm, subcolspace=8mm]{tikzposter}

\usepackage[utf8]{inputenc}
\usepackage{graphicx}
\usepackage{booktabs}
\usepackage{qrcode}
\usepackage{enumitem}
\setlist[itemize]{leftmargin=1.2em, itemsep=0.3em, topsep=0.3em}

\usetheme{Simple}
\usecolorstyle{Denmark}
\colorlet{backgroundcolor}{white}
\colorlet{framecolor}{titlebgcolor}

\newcommand{\fig}[3][1.0]{%
  \begin{tikzfigure}\includegraphics[width=#1\linewidth]{#2}\end{tikzfigure}
  \textit{\small #3}}

\title{\parbox{0.92\linewidth}{\centering
  Beyond Speed Rankings: A Nine-Country Measurement Study of\\
  Southeast Asia's Application-Layer Digital Divide}}
\author{Chissanupun Athiwarikanon, Kunanont Malayanont, Pakkapon Pattanakul}
\institute{Kasetsart University}

\begin{document}
\maketitle

% ---- Headline stat band ----------------------------------------------
\begin{columns}
  \column{0.25}
  \block{}{\centering \color{titlebgcolor}\Huge\textbf{990M}\\[0.3em]
    \normalsize\color{black} Speed tests analyzed}
  \column{0.25}
  \block{}{\centering \color{titlebgcolor}\Huge\textbf{9}\\[0.3em]
    \normalsize\color{black} SE Asian countries}
  \column{0.25}
  \block{}{\centering \color{titlebgcolor}\Huge\textbf{100\%}\\[0.3em]
    \normalsize\color{black} HD-video success, every country}
  \column{0.25}
  \block{}{\centering \color{titlebgcolor}\Huge\textbf{$<$38\%}\\[0.3em]
    \normalsize\color{black} Cloud-gaming success, weakest markets}
\end{columns}

\begin{columns}

  % ==================================================================
  \column{0.33}
  \block{Motivation}{
    Speed-index rankings say little about whether real applications work.
    Thailand ranked \textbf{12th worldwide} for fixed broadband (Ookla
    Speedtest Global Index, June 2026), yet a 2023 Thailand Consumer
    Council survey found \textbf{81\%} of 2{,}924 respondents had
    connection problems --- mostly slow and dropping connections. The
    mismatch is not unique to Thailand.

    \vspace{0.8em}
    \textbf{Question:} how do nine Southeast Asian networks perform against
    the demands of real applications --- voice, HD video, UHD video and
    cloud gaming?
  }

  \block{Data}{
    \begin{itemize}
      \item \textbf{Ookla Open Data:} 229.4M tests, quarterly ${\sim}$600\,m
        tiles, fixed + mobile
      \item \textbf{M-Lab NDT7:} 760.6M single-stream tests
      \item \textbf{Coverage:} Thailand, Vietnam, Laos, Myanmar, Cambodia,
        Indonesia, Philippines, Malaysia, Singapore; Q1 2023--Q4 2025
      \item Tiles mapped to provinces (point-in-polygon, geoBoundaries ADM1),
        test-weighted; province-quarters with $\geq$100 tests and $\geq$5 tiles
    \end{itemize}
  }

  \block{Where the Speed Is}{
    The five fastest provinces per country, hatched where a province carries
    under \textbf{5\%} of its country's tests and its rank is therefore not
    reliable. Singapore clears \textbf{360~Mbps} in every region; Thailand's
    top five are all Bangkok-adjacent and all thin-sample. Malaysia and the
    Philippines concentrate in hubs that do carry real volume (Selangor
    34.3\%, NCR 28.0\%), while the emerging markets top out below 200~Mbps
    and fill their rankings with low-sample provinces.

    \vspace{0.5em}
    \fig[0.95]{../paper/figures/image5.png}{Top five provinces per country, fixed broadband; hatched = under 5\% of national tests.}
  }

  \block{Data \& Code}{
    \begin{center}
      \qrcode[height=5cm]{https://github.com/chissanupun/thailand-broadband-analysis}
      \\[0.6em] \normalsize \texttt{github.com/chissanupun/}\\
      \texttt{thailand-broadband-analysis}
    \end{center}
  }

  % ==================================================================
  \column{0.34}
  \block{Finding 1 --- Basics Are Solved}{
    Voice ($\geq$64~kbps) and HD video ($\geq$5~Mbps) succeed at
    \textbf{100\%} in every country, every quarter, on fixed and mobile
    networks. Even the lowest-tier broadband markets meet the basic floor.

    \vspace{0.6em}
    \begin{center}
    \begin{tabular}{lll}
      \toprule
      \textbf{Application} & \textbf{Data rate} & \textbf{Latency} \\
      \midrule
      Voice        & 64 kbps & 200 ms \\
      Video (HD)   & 5 Mbps  & few s \\
      Video (UHD)  & 25 Mbps & few s \\
      Cloud gaming & 44 Mbps & 25 ms \\
      \bottomrule
    \end{tabular}
    \end{center}
    \small Thresholds after L\"ubben \& Misfeld (2022). Success rate =
    national-level share meeting the download-speed threshold, per quarter.
  }

  \block{Finding 2 --- Urban--Rural Stratification}{
    Most capitals outperform their nation: Bangkok 298.8 vs.\ 276.9~Mbps,
    Kuala Lumpur 212.3 vs.\ 194.8, Manila NCR 159.6 vs.\ 132.0.
    Singapore's Central Region (361.0) sits \emph{below} its national mean
    (394.8). Peak capital speed is \textbf{over 16$\times$} the
    lowest capital (Myanmar, 22.0~Mbps).

    \vspace{0.5em}
    \fig{../../outputs/ookla/cross_country/03_capital_vs_national.png}{Capital province vs.\ national mean, fixed broadband, test-weighted.}
  }

  \block{Finding 3 --- Speed Gap Is Widening}{
    Singapore nearly \textbf{doubled} its median provincial speed
    (${\sim}$280 $\to$ ${\sim}$550~Mbps); Thailand grew ${\sim}$200 $\to$
    ${\sim}$300. Vietnam accelerated to catch Malaysia (200+~Mbps).
    Laos, Cambodia, Indonesia and Myanmar stayed \textbf{below 60~Mbps}
    for all three years.

    \vspace{0.5em}
    \fig{../paper/figures/image6.png}{Median provincial fixed download speed, Q1 2023--Q4 2025.}
  }

  % ==================================================================
  \column{0.33}
  \block{Finding 4 --- Application Tiers Diverge}{
    \textbf{UHD video ($\geq$25~Mbps)}
    \begin{itemize}
      \item Malaysia, Singapore, Thailand, Philippines, Vietnam: 100\% every quarter
      \item Cambodia 67--86\%; Indonesia 36 $\to$ 65\%; Myanmar 40--55\%
      \item Laos volatile: 50--100\%
    \end{itemize}
    \fig[0.82]{../../outputs/app_success_rate_median/fig_video_uhd_national.png}{UHD-video success rate by country.}

    \vspace{0.8em}
    \textbf{Cloud gaming ($\geq$44~Mbps)}
    \begin{itemize}
      \item Singapore ${\sim}$100\%; Thailand and Malaysia 87--100\%
      \item Philippines 58--93\%; Vietnam $<$50\% $\to$ ${\sim}$94\%
      \item Cambodia 38--49\%; Laos ${\sim}$52\%
      \item Indonesia $<$38\%; Myanmar $<$31\%
    \end{itemize}
    \fig[0.82]{../../outputs/app_success_rate_median/fig_cloud_gaming_national.png}{Cloud-gaming success rate by country.}
  }

  \block{Finding 5 --- Mobile Is Uniform, Fixed Is Not}{
    Fixed broadband beats mobile in most countries: Thailand shows the
    largest gap (\textbf{3.2$\times$}, ${\sim}$275 vs.\ ${\sim}$85~Mbps),
    Singapore 2.1$\times$, the Philippines 1.9$\times$. Malaysia and
    Indonesia are near parity (1.1$\times$). Myanmar is the only country
    where mobile is \emph{faster} than fixed (0.8$\times$) --- a
    mobile-first market.

    \vspace{0.5em}
    \fig[0.92]{../paper/figures/image4.png}{Fixed vs.\ mobile mean download by country, with ratio.}
  }

  \block{Takeaway}{
    Raw speed rankings hide the divide that matters: not whether a country
    ``has broadband,'' but whether it survives a real application. Basic
    needs are met everywhere; UHD and cloud gaming expose the gap. This
    argues for targeted fiber deployment, regional investment and
    transport-layer optimization.
  }

\end{columns}

\end{document}
